\section*{Capítulo \ref{ch:g10:measuring-the-earth} --- Medir la Tierra}

\begin{solution}{exo:g10:measuring-the-earth:1}
Un barco desaparece empezando por el casco: sobre un mar plano solo se
haría más pequeño y seguiría viéndose entero. La sombra de la Tierra sobre la
Luna durante un eclipse es siempre un arco de circunferencia: un disco plano
visto de canto proyectaría una sombra alargada. Al viajar hacia el norte, la
estrella Polar sube sobre el horizonte: en una Tierra plana su altura apenas
cambiaría.
\end{solution}

\begin{solution}{exo:g10:measuring-the-earth:2}
La \omterm{def:g10:measuring-the-earth:geodesy}{geodesia} mide la forma y el tamaño de la Tierra y la posición de los
puntos de su superficie, y cómo cambian con el tiempo. La \omterm{def:g10:measuring-the-earth:triangulation}{triangulación}
localiza puntos mediante una cadena de triángulos en la que se miden todos
los ángulos; solo una longitud, la base, se mide sobre el terreno.
\end{solution}

\begin{solution}{exo:g10:measuring-the-earth:3}
$C = 50 \times 5000 = 250\,000$ estadios; $250\,000 \times \qty{157.2}{m} =
\qty{39300}{km}$.
\end{solution}

\begin{solution}{exo:g10:measuring-the-earth:4}
$111/60 \approx \qty{1.85}{km}$. Cuarto de meridiano: $90 \times 111 \approx
\qty{10000}{km}$.
\end{solution}

\begin{solution}{exo:g10:measuring-the-earth:5}
$a - b = 6378.137 - 6356.752 = \qty{21.385}{km}$;
$f = 21.385/6378.137 \approx 0.003353 \approx 1/298$.
\end{solution}

\begin{solution}{exo:g10:measuring-the-earth:6}
Sur: $\tan\theta_1 = 0.364$, $\theta_1 \approx \qty{20.0}{\degree}$. Norte:
$\tan\theta_2 = 0.456$, $\theta_2 \approx \qty{24.5}{\degree}$.
$\alpha = \qty{4.5}{\degree}$, luego $C = (360/4.5) \times 500 = 80 \times 500 =
\qty{40000}{km}$.
\end{solution}

\begin{solution}{exo:g10:measuring-the-earth:7}
$C = 180 - 62 - 71 = \qty{47}{\degree}$.
$AC = 10.00 \sin 71^\circ/\sin 47^\circ = 10.00 \times 0.9455/0.7314 \approx
\qty{12.93}{km}$; $BC = 10.00 \sin 62^\circ/\sin 47^\circ \approx
\qty{12.07}{km}$.
\end{solution}

\begin{solution}{exo:g10:measuring-the-earth:8}
El cambio de longitud es $10 + 20 = \qty{30}{\degree}$. A
\qty{60}{\degree} N un grado mide $111.3 \times 0.5 = \qty{55.65}{km}$, así que
la distancia es $30 \times 55.65 \approx \qty{1670}{km}$. Sobre el ecuador:
$30 \times 111.3 \approx \qty{3340}{km}$, el doble.
\end{solution}

\begin{solution}{exo:g10:measuring-the-earth:9}
La visual toca el mar en $T$: la tangente a una circunferencia es
perpendicular al radio en el punto de contacto, así que el ángulo en $T$ es
recto. Por Pitágoras, $(R+h)^2 = R^2 + d^2$, luego $d^2 = 2Rh + h^2$; $h$ es
diminuto frente a $R$, así que $d \approx \sqrt{2Rh} = \sqrt{2 \times
6\,371\,000 \times 2} \approx \qty{5050}{m}$, unos \qty{5}{km}.
\end{solution}

\begin{solution}{exo:g10:measuring-the-earth:10}
La curva es más empinada en las latitudes medias, de unos \qty{30}{\degree} a
\qty{60}{\degree}, y casi plana cerca del ecuador y del polo. En una Tierra
achatada, la superficie está menos curvada cerca de los polos; cuanto menos
curvada está la superficie, más hay que recorrer para que la vertical (y, con
ella, la latitud) gire un grado.
\end{solution}

\begin{solution}{exo:g10:measuring-the-earth:11}
$1072.2/9.646 \approx \qty{111.15}{km}$ por grado; $90 \times 111.15
\approx \qty{10004}{km}$, \qty{2}{km} más que el verdadero cuarto de
meridiano. El arco está centrado cerca de \qty{46}{\degree}, donde un grado
es algo más largo que el grado medio del ecuador al polo
($10\,002.0/90 \approx \qty{111.13}{km}$).
\end{solution}

\begin{solution}{exo:g10:measuring-the-earth:12}
Con \qty{7.1}{\degree}: $C = (360/7.1) \times 5000 \approx 253\,500$ estadios
en lugar de $250\,000$: alrededor de un 1.4\,\% de más (un ángulo menor para
la misma distancia significa una Tierra mayor). La distancia, medida a pasos
por caminantes a lo largo del Nilo, era mucho menos precisa que el ángulo.
\end{solution}

\begin{solution}{exo:g10:measuring-the-earth:13}
En $ABC$: $C = \qty{50}{\degree}$, $AC = 6.000 \sin 75^\circ/\sin 50^\circ
\approx \qty{7.566}{km}$. En $ACD$: el ángulo en $D$ es $180 - 50 - 68 =
\qty{62}{\degree}$, y $AD = 7.566 \sin 68^\circ/\sin 62^\circ \approx
\qty{7.945}{km}$.
\end{solution}

\begin{solution}{exo:g10:measuring-the-earth:14}
$\cos\varphi = 1/2$ da $\varphi = \qty{60}{\degree}$; $\cos\varphi = 0.1$
da $\varphi \approx \qty{84.3}{\degree}$. Un grado de longitud a la latitud
$\varphi$ es el grado del ecuador multiplicado por $\cos\varphi$: la razón
depende solo del ángulo, no del radio.
\end{solution}

\begin{solution}{exo:g10:measuring-the-earth:15}
En el ecuador: $111.195 - 110.574 = \qty{0.621}{km} = \qty{621}{m}$; en el
polo: $111.694 - 111.195 = \qty{0.499}{km} = \qty{499}{m}$. Errores de varios
cientos de metros por grado, que se acumulan grado tras grado, son muchísimo
mayores que el metro: las posiciones calculadas a partir de latitudes y
longitudes tienen que usar el elipsoide.
\end{solution}

\begin{solution}{pb:g10:measuring-the-earth:1}
\textbf{1.} $31.20 - 24.09 = \qty{7.11}{\degree}$.

\textbf{2.} $7.11/360 \approx 0.0198 = 1/50.6$, muy cerca del $1/50$ de
Eratóstenes.

\textbf{3.} $C = 788 \times 360/7.11 \approx \qty{39900}{km}$.

\textbf{4.} El arco de meridiano entre los paralelos: la diferencia entre
los ángulos de mediodía es la diferencia de latitud, que solo mide la
separación norte--sur de los dos lugares; la componente este--oeste de la
distancia no interviene.

\textbf{5.} A mediodía local el Sol está en el plano del meridiano, el mismo
para los dos lugares; el mismo día ocupa la misma posición entre las
estrellas, de modo que los dos ángulos solo difieren por la curvatura de la
Tierra.

\textbf{6.} $180 - 58 - 66 = \qty{56}{\degree}$.

\textbf{7.} $AC = 12.000 \sin 66^\circ/\sin 56^\circ = 12.000 \times
0.9135/0.8290 \approx \qty{13.223}{km}$.

\textbf{8.} $BC = 12.000 \sin 58^\circ/\sin 56^\circ \approx \qty{12.275}{km}$.

\textbf{9.} El ángulo en $D$ es $180 - 47 - 75 = \qty{58}{\degree}$;
$AD = 13.223 \sin 75^\circ/\sin 58^\circ \approx \qty{15.061}{km}$.

\textbf{10.} $15.061 \cos 8^\circ \approx \qty{14.914}{km}$.

\textbf{11.} $14.914/111.13 \approx \qty{0.1342}{\degree}$, es decir,
$0.1342 \times 60 \approx 8.05$ minutos de arco.

\textbf{12.} $111.51 - 110.57 = \qty{0.94}{km}$: \qty{940}{m}.

\textbf{13.} Cerca del polo, la vertical gira más despacio con la distancia:
allí la superficie está menos curvada. La Tierra está achatada por los polos
y abultada en el ecuador.

\textbf{14.} $a - b = \qty{21.385}{km}$; $f = 21.385/6378.137 \approx 1/298$.

\textbf{15.} La plomada es perpendicular a la superficie del elipsoide. En
una elipse, la perpendicular a la curva solo pasa por el centro en el ecuador
y en los polos; a \qty{45}{\degree} corta el eje por debajo del centro.

\textbf{16.} $10\,000\,000$ metros, es decir, \qty{10000}{km}.

\textbf{17.} $10\,001\,966/10\,000\,000 \approx \qty{1.0001966}{m}$.

\textbf{18.} Demasiado corto: la verdadera diezmillonésima parte es
$1.0001966$ veces el metro de 1799, así que al metro de 1799 le falta
alrededor de $1.97 \times 10^{-4}$ de su longitud, un 0.02\,\%.

\textbf{19.} Demasiado corto. La barra es la diezmillonésima parte del
cuarto de meridiano tal como se calculó en la década de 1790; como la barra
es más corta que la diezmillonésima parte del verdadero cuarto de meridiano,
el cuarto de meridiano calculado era más corto que el verdadero.

\textbf{20.} $1.0001966 - 1 = 0.0001966$ m: el primer metro era unos
\textbf{\qty{0.2}{mm} demasiado corto}, el grosor de unas dos hojas de papel
en un metro, una precisión asombrosa para una medición hecha con anteojos y
ángulos a través de un país entero.
\end{solution}
